\documentclass[12pt,a4paper]{article}
\usepackage[utf8]{inputenc}
\usepackage[italian]{babel}
\usepackage{Alegreya}
\usepackage[T1]{fontenc}
\usepackage{geometry}
\geometry{margin=2.5cm}
\usepackage{titlesec}
\usepackage{xcolor}
\usepackage{mdframed}

% Stile per i titoli
\titleformat{\section}{\Large\bfseries\color{darkgray}}{\thesection.}{1em}{}
\titleformat{\subsection}{\large\bfseries\color{gray}}{\thesubsection}{1em}{}

% Ambiente per gli estratti
\newmdenv[
  backgroundcolor=gray!5,
  linecolor=gray!30,
  linewidth=1pt,
  leftmargin=1cm,
  rightmargin=1cm,
  innertopmargin=10pt,
  innerbottommargin=10pt
]{estratto}

% Ambiente per le riflessioni
\newmdenv[
  backgroundcolor=blue!2,
  linecolor=blue!20,
  linewidth=1pt,
  leftmargin=0.5cm,
  rightmargin=0.5cm,
  innertopmargin=8pt,
  innerbottommargin=8pt
]{riflessione}

\title{\textbf{Filo del discorso:}\\ \large Utopia-Anacronismo negli scritti di Mario Bertoncini}
\author{Appunti per l'introduzione all'ascolto\\ \textit{LAZZARO IN RIPRESA DI VENTI}}
\date{}

\begin{document}
\maketitle

\section{La condizione dell'arte nel presente: essere fuori tempo}

\begin{estratto}
\textbf{Estratto:} \\
\textit{``Molto s'è detto e scritto, in tutti i tempi ma soprattutto nel nostro, dei limiti, degli scopi, dei destini e persino della morte dell'arte e senza dubbio, a dispetto di quanto detto o scritto con maggiore o minore autorevolezza, d'arte si continua a parlare e spesso a torto o a ragione, soprattutto a farne. Non sta a noi, che d'arte e per l'arte viviamo, decidere quando essa sia legittima, quando cioè [...] essa appaia in sintonia con le esigenze della società circostante: anche se esattamente questo sembrerebbe mancare oggi in un mondo oppresso da interrogativi irrisolti, da tragedie ricorrenti e dalle angosce d'un quotidiano privo di adeguata soluzione.''}
\end{estratto}

\begin{riflessione}
\textbf{Riflessione:} \\
Bertoncini parte dalla condizione anacronistica dell'arte contemporanea: essa persiste \textit{nonostante} la sua apparente illegittimità rispetto al tempo presente. Non è ``in sintonia'' con la società, eppure continua a esistere. Questa discronia non è difetto ma condizione essenziale: l'arte è anacronistica perché deve esserlo per cogliere ciò che il presente nasconde a se stesso. È il primo movimento del ``contemporaneo'' agambeniano: la non-coincidenza con il proprio tempo.

\vspace{0.5cm}
\textit{[Spazio per appunti e integrazioni]}
\vspace{1cm}
\end{riflessione}

\section{L'alternativa utopica: dal suono immaginario allo spazio del vento}

\begin{estratto}
\textbf{Estratto:} \\
\textit{``Quale può dunque essere l'alternativa desiderabile per un'arte amata oppure odiata ma comunque a dispetto di tutto sempre presente tuttavia in un'epoca sia pure distratta e turbolenta come quella nostra? O meglio, per scoprire le carte, qual è la soluzione che noi caldeggiamo, volendo portare il suono del nostro immaginario a confronto con uno spazio più vasto, come è quello percorso dal vento o dalla magia delle campane?''}
\end{estratto}

\begin{riflessione}
\textbf{Riflessione:} \\
Qui emerge la dimensione utopica: non come evasione ma come espansione dell'``immaginario'' verso ``uno spazio più vasto''. Il vento e le campane sono elementi che attraversano e connettono spazi e tempi diversi - sono agenti di anacronismo sonoro. L'utopia bertonciana non è un altrove ma un'intensificazione del qui, un'apertura verso il possibile nascosto nel reale.

\vspace{0.5cm}
\textit{[Spazio per appunti e integrazioni]}
\vspace{1cm}
\end{riflessione}

\section{La pratica archeologica: scavare nel presente per trovare l'inatteso}

\begin{estratto}
\textbf{Estratto:} \\
\textit{``Nel tipo particolare di musica che si vuole descrivere qui, suonare, eseguire, è un'azione idealmente e metaforicamente assimilabile a quella di scavo. Come l'archeologo libera accuratamente dalle incrostazioni e dai detriti che lo ricoprono il reperto ambito, cercato ma non conosciuto in anticipo, così all'esecutore attento si chiede in questo caso di scegliere fra i tanti percorsi possibili lungo la densa compagine di corde quell'unico gesto o quella serie di gesti già esplicitamente catalogati nelle note per l'esecuzione.''}
\end{estratto}

\begin{riflessione}
\textbf{Riflessione:} \\
Ecco il metodo anacronistico per eccellenza: lo scavo. L'esecutore non riproduce un passato noto ma cerca ``il reperto ambito, cercato ma non conosciuto in anticipo''. È l'operazione del contemporaneo che scava nelle ``tenebre'' del presente per trovare una luce che non sapeva di cercare. L'anacronismo qui diventa tecnica compositiva: ogni esecuzione è un'archeologia del possibile.

\vspace{0.5cm}
\textit{[Spazio per appunti e integrazioni]}
\vspace{1cm}
\end{riflessione}

\section{L'interpretazione multipla: liberarsi dall'univocità temporale}

\begin{estratto}
\textbf{Estratto:} \\
\textit{``L'insieme delle misure necessarie alla realizzazione, allo sviluppo sonoro dell'oggetto, sono indispensabili per l'esecutore che voglia tradurre in atto una interpretazione adeguata di esso, ma non sono da confondere con un piano univoco di sviluppo temporale. [...] tutte le volte che io stesso definisco graficamente (o verbalmente) una determinata partitura, non detto in modo univoco l'unica interpretazione dell'oggetto, ma soltanto una delle molteplici interpretazioni di esso e sono quindi uno dei tanti possibili interpreti della mia stessa composizione.''}
\end{estratto}

\begin{riflessione}
\textbf{Riflessione:} \\
Bertoncini rompe con la temporalità lineare e univoca della tradizione musicale occidentale. Il compositore stesso diventa ``uno dei tanti possibili interpreti'' - posizione perfettamente anacronistica che dissolve la gerarchia temporale tra creazione, interpretazione e fruizione. Ogni momento esecutivo può essere contemporaneo a se stesso in modi diversi.

\vspace{0.5cm}
\textit{[Spazio per appunti e integrazioni]}
\vspace{1cm}
\end{riflessione}

\section{L'objet trouvé onirico: quando l'anacronismo si fa sogno}

\begin{estratto}
\textbf{Estratto:} \\
\textit{``l'osservatore–ricettore–esecutore dell'oggetto, nell'atto cosciente dell'applicare alla scultura di suono una successione qualsiasi dei modelli esecutivi memorizzati, nel rivelarlo scopre inconsciamente un ductus sonoro agente sulla propria immaginazione; e la sorpresa che ne deriva agisce su quella con la stessa vivida sensazione che accompagna il sogno. E come nel sogno, anche nel processo eolico di lettura, in quel groviglio di corde delle slegate immagini oniriche si ricollegano continuamente ad una rete di ricordi allusivi la cui assurdità fantastica viene proposta in modo momentaneamente credibile.''}
\end{estratto}

\begin{riflessione}
\textbf{Riflessione:} \\
Il sogno come modello di temporalità anacronistica: nel sogno passato, presente e futuro si ricompongono in configurazioni impossibili ma ``momentaneamente credibili''. L'``objet trouvé'' è ciò che affiora da questo groviglio temporale - non è un ritrovamento archeologico ma una creazione onirica. L'anacronismo qui si fa poetico, generativo di senso inatteso.

\vspace{0.5cm}
\textit{[Spazio per appunti e integrazioni]}
\vspace{1cm}
\end{riflessione}

\section{L'utopia radicale: ``perché limitarne la portata per amore d'un banale realismo?''}

\begin{estratto}
\textbf{Estratto:} \\
\textit{``MAESTRO: Beh, ero sul punto di dire che quel progetto campato in aria, anziché ricalcare utopicamente l'ambiziosa idea originale delle tre grandi arpe eolie da installare in montagna, all'aperto — un'idea che dati i costi sono riuscito poi a eseguire soltanto in forma ridotta —, nella leggerezza del sogno si presentò in una veste ancora più impegnativa e dispendiosa. Dato che per me quell'idea era irrealizzabile, perché limitarne la portata per amore d'un banale realismo?''}
\end{estratto}

\begin{riflessione}
\textbf{Riflessione:} \\
Qui l'utopia si rivela nella sua essenza paradossale: più è impossibile, più diventa necessaria. Il ``progetto campato in aria'' non è fuga dal reale ma sua intensificazione onirica. Bertoncini compie il gesto utopico per eccellenza: di fronte all'impossibilità, non riduce ma amplifica il desiderio. È l'anacronismo che si fa programmatico: sognare oltre ogni realismo per toccare il possibile.

\vspace{0.5cm}
\textit{[Spazio per appunti e integrazioni]}
\vspace{1cm}
\end{riflessione}

\section{Il messaggero ermetico: unificare reale e ideale}

\begin{estratto}
\textbf{Estratto:} \\
\textit{``L'interpretazione è la volontà che il mondo abbia un senso. Hérmes, il messaggero. Unifica, attiva, le due dimensioni: reale e ideale.''}
\end{estratto}

\begin{riflessione}
\textbf{Riflessione:} \\
Citazione di Severino che chiude il cerchio: Hermes è la figura dell'anacronismo per eccellenza, colui che attraversa i tempi e i mondi portando messaggi tra dimensioni diverse. L'``interpretazione'' musicale è questo gesto ermetico che ``unifica'' reale e ideale non per sintesi dialettica ma per attraversamento, per ``attivazione'' di connessioni impensate.

\vspace{0.5cm}
\textit{[Spazio per appunti e integrazioni]}
\vspace{1cm}
\end{riflessione}

\newpage

\section*{Sintesi del filo discorsivo}

Il percorso negli scritti di Bertoncini rivela una concezione dell'arte musicale come pratica essenzialmente anacronistica, che trova nell'utopia non un rifugio ma il proprio metodo di lavoro. L'anacronismo qui non è nostalgia ma tecnica di scavo nel presente; l'utopia non è evasione ma radicalizzazione dell'immaginazione.

Il dialogo tra \textsc{Bertoncini} e \textsc{Lazzaro} può articolarsi lungo questo filo: il compositore che scava nelle tenebre sonore del presente per far emergere oggetti sonori impensati, e l'interprete-resuscitato che dà corpo a queste visioni attraverso il soffio, il vento, la ``ripresa'' vitale. Entrambi praticano quella forma di contemporaneità che sa essere fuori tempo per cogliere il tempo nella sua verità nascosta.

\vspace{2cm}

\section*{Note aggiuntive per il concerto}

\vspace{2cm}
\textit{[Spazio per sviluppi del dialogo BERTONCINI-LAZZARO]}

\vspace{5cm}

\section*{Riferimenti}

\begin{itemize}
\item Giorgio Agamben, \textit{Che cos'è il contemporaneo?}
\item Mario Bertoncini, \textit{Ragionamenti in forma di dialogo}
\item Emanuele Severino, \textit{La filosofia futura}
\item Ingeborg Bachmann, \textit{Letteratura come Utopia}
\end{itemize}

\end{document}
